% Options for packages loaded elsewhere
\PassOptionsToPackage{unicode}{hyperref}
\PassOptionsToPackage{hyphens}{url}
%
\documentclass[10pt,twocolumn]{article}
\usepackage{amsmath,amssymb}
\usepackage[letterpaper,top=0.66in,bottom=0.66in,left=0.66in,right=0.66in,columnsep=0.24in]{geometry}
\usepackage{tcolorbox}
\usepackage{dblfloatfix}
\usepackage{caption}
\usepackage{float}
\usepackage{tabularx}
\usepackage{array}
\usepackage{enumitem}
\usepackage{fancyhdr}
\usepackage{titlesec}
\usepackage{iftex}
\ifPDFTeX
  \usepackage[T1]{fontenc}
  \usepackage[utf8]{inputenc}
  \usepackage{textcomp} % provide euro and other symbols
\else % if luatex or xetex
  \usepackage{unicode-math} % this also loads fontspec
  \defaultfontfeatures{Scale=MatchLowercase}
  \defaultfontfeatures[\rmfamily]{Ligatures=TeX,Scale=1}
\fi
\usepackage{lmodern}
\ifPDFTeX\else
  % xetex/luatex font selection
\fi
% Use upquote if available, for straight quotes in verbatim environments
\IfFileExists{upquote.sty}{\usepackage{upquote}}{}
\IfFileExists{microtype.sty}{% use microtype if available
  \usepackage[]{microtype}
  \UseMicrotypeSet[protrusion]{basicmath} % disable protrusion for tt fonts
}{}
\setlength{\parindent}{1em}
\setlength{\parskip}{0pt}
\usepackage{xcolor}
\usepackage{longtable,booktabs,array}
\usepackage{calc} % for calculating minipage widths
% Correct order of tables after \paragraph or \subparagraph
\usepackage{etoolbox}
\makeatletter
\patchcmd\longtable{\par}{\if@noskipsec\mbox{}\fi\par}{}{}
\makeatother
% Allow footnotes in longtable head/foot
\IfFileExists{footnotehyper.sty}{\usepackage{footnotehyper}}{\usepackage{footnote}}
\makesavenoteenv{longtable}
\usepackage{graphicx}
\makeatletter
\def\maxwidth{\ifdim\Gin@nat@width>\linewidth\linewidth\else\Gin@nat@width\fi}
\def\maxheight{\ifdim\Gin@nat@height>\textheight\textheight\else\Gin@nat@height\fi}
\makeatother
% Scale images if necessary, so that they will not overflow the page
% margins by default, and it is still possible to overwrite the defaults
% using explicit options in \includegraphics[width, height, ...]{}
\setkeys{Gin}{width=\maxwidth,height=\maxheight,keepaspectratio}
% Set default figure placement to htbp
\makeatletter
\def\fps@figure{htbp}
\makeatother
\setlength{\emergencystretch}{3em} % prevent overfull lines
\providecommand{\tightlist}{%
  \setlength{\itemsep}{0pt}\setlength{\parskip}{0pt}}
\setcounter{secnumdepth}{-\maxdimen}
\setlength{\columnsep}{0.24in}
\setlength{\floatsep}{8pt plus 2pt minus 2pt}
\setlength{\textfloatsep}{9pt plus 2pt minus 2pt}
\setlength{\intextsep}{8pt plus 2pt minus 2pt}
\renewcommand{\topfraction}{0.92}
\renewcommand{\dbltopfraction}{0.92}
\renewcommand{\textfraction}{0.06}
\renewcommand{\floatpagefraction}{0.75}
\renewcommand{\dblfloatpagefraction}{0.75}
\titlespacing*{\section}{0pt}{1.05em}{0.45em}
\titlespacing*{\subsection}{0pt}{0.8em}{0.3em}
\pagestyle{fancy}
\fancyhf{}
\fancyhead[L]{\small RobustLens}
\fancyhead[R]{\small Compound-corruption robustness}
\fancyfoot[C]{\thepage}
\renewcommand{\headrulewidth}{0.35pt}
\setlength{\headheight}{13pt}
\captionsetup{font=small,labelfont=bf,justification=centering}
\tcbset{colback=black!94,colframe=black!94,coltext=white,boxrule=0.5pt,arc=1.5mm,left=2.5mm,right=2.5mm,top=1.5mm,bottom=1.5mm} % remove section numbering
\ifLuaTeX
  \usepackage{selnolig}  % disable illegal ligatures
\fi
\usepackage{bookmark}
\IfFileExists{xurl.sty}{\usepackage{xurl}}{} % add URL line breaks if available
\urlstyle{same}
\hypersetup{
  pdftitle={Order Matters Beyond Accuracy: Paired Prediction Analysis of Compound Image Corruptions},
  pdfauthor={Suryakanth V. Gangashetty; Aviral Choudhary},
  hidelinks,
  pdfcreator={LaTeX via pandoc}}

\title{Order Matters Beyond Accuracy: Paired Prediction Analysis of
Compound Image Corruptions}
\author{Suryakanth V. Gangashetty$^{1}$\\[3pt] Aviral Choudhary$^{2}$}
\date{October 2026}

\begin{document}
\twocolumn[
\begin{@twocolumnfalse}
\vspace*{-0.45em}
\begin{center}
{\LARGE\bfseries Order Matters Beyond Accuracy: Paired Prediction Analysis of Compound Image Corruptions\par}
\vspace{0.65em}
{\large Suryakanth V. Gangashetty$^{1}$\par}
\vspace{0.14em}
{\large Aviral Choudhary$^{2}$\par}
\vspace{0.24em}
{\footnotesize $^{1}$Department of Computer Science and Engineering, Koneru Lakshmaiah Education Foundation, India\par}
{\footnotesize $^{2}$Independent Researcher, India\par}
\vspace{0.32em}
{\small October 2026\par}
\end{center}
\vspace{0.45em}
{\small\noindent\textbf{Abstract.} Classification accuracy under compound corruption depends on the order
in which degradations are applied. We evaluate all six permutations of
Gaussian noise (N), Gaussian blur (B), and JPEG compression on 3,925
Imagenette-v2 validation images using four pretrained classifiers:
ResNet-50, ConvNeXt-Tiny, ViT-Base/16, and Swin-Tiny. Under common
nominal settings, the model-specific six-order accuracy ranges span
17.13 to 21.55 percentage points. Two orders that both end in JPEG
differ by about 16 percentage points in across-model mean accuracy. A
second parameter regime selected toward shared PSNR and SSIM targets
reduces the range of the six across-model mean accuracies from 18.43 to
6.79 percentage points, a 63.18\% descriptive contraction. A paired
five-state analysis shows that predicted-label churn is 40.78\% under
nominal settings and 34.74\% under the PSNR/SSIM-targeted regime,
compared with correctness-state churn of 16.93\% and 12.34\%. Many of
the additional label changes occur between two incorrect classes,
changes that an accuracy-only analysis cannot show.
\par\vspace{0.35em}}
{\small\noindent\textbf{Keywords:} computer vision; corruption robustness; compound
degradations; operator order; paired prediction analysis; image fidelity}
\vspace{0.6em}
\begin{tcolorbox}
{\large\bfseries WHAT THIS STUDY FINDS\par}\vspace{0.25em}
\small
\begin{itemize}
\tightlist
\item
  \textbf{Processing order and accuracy.} The model-specific six-order
  accuracy ranges span 17.13--21.55 percentage points. NBJ and BNJ both
  end in JPEG but differ by about 16 percentage points in across-model
  mean accuracy.
\item
  \textbf{PSNR/SSIM-targeted settings narrow the aggregate range.} The
  range across the six across-model mean accuracies decreases from 18.43
  to 6.79 percentage points (63.18\%). Because the prespecified
  PSNR/SSIM tolerances were not met, this is a descriptive comparison of
  the two tested parameter regimes.
\item
  \textbf{Predicted labels change beyond correctness.} Nominal
  predicted-label churn is 40.78\%, compared with 16.93\%
  correctness-state churn. A further 23.85\% of paired observations
  remain incorrect while switching to a different incorrect class.
\end{itemize}
\end{tcolorbox}
\end{@twocolumnfalse}
]

\section{1. Introduction}\label{introduction}

Images used by vision systems often undergo several operations during
sensing, processing, compression, or transmission. Gaussian noise, blur,
and JPEG compression are familiar degradations, but applying the same
operations in a different order can change the resulting image and
classifier output. Standard robustness benchmarks evaluate many
degradations in isolation; work on dual corruptions, image-processing
chains, and compositional robustness motivates evaluating composed
pipelines {[}1-3, 6, 9{]}.

This study asks: \textbf{when the same source images and fixed model
weights are exposed to the same three degradations in different orders,
how much do accuracy and individual predicted labels change?} We
evaluate all six permutations of Gaussian noise, Gaussian blur, and JPEG
compression on 3,925 Imagenette-v2 validation images using four
pretrained classifiers. We compare common nominal parameters with a
second regime in which sequence-specific settings are selected toward
shared PSNR and SSIM targets.

For each order, we report top-1 accuracy and image-fidelity summaries,
then compare predicted classes for the same image-severity observations
across order pairs. This separates changes in correctness from switches
between incorrect labels. The five-state partition is used as an
analysis tool; the contribution is its application to the complete
six-order experiment, alongside the accuracy and fidelity comparisons.

\section{2. Related work}\label{related-work}

ImageNet-C and ImageNet-P established widely used protocols for
corruption robustness and perturbation stability {[}1{]}. Goyal et
al.~studied noise-blur dual corruptions {[}2{]}, and Liao et
al.~examined image operator-chain detection {[}3{]}. Other work has
studied compositional robustness {[}9{]}, real-world corruption chains
for large multimodal models {[}5{]}, continuous nuisance shifts {[}6{]},
and out-of-distribution evaluation of web-scale vision models {[}7{]}.
Prediction churn has also been studied across model versions {[}4{]}.

The closest methodological context is recent paired-evaluation work
under visual perturbations {[}8{]}, which discusses a five-state
refinement that distinguishes persistent wrong answers from changed
wrong answers. The five-state partition in this paper is therefore not
presented as a new general method. This study applies paired prediction
comparisons to all six permutations of one fixed noise-blur-JPEG chain
across four frozen classifiers and reports the results alongside
order-specific accuracy and image-fidelity summaries.

\section{3. Experimental design}\label{experimental-design}

\subsection{3.1 Dataset and
classifiers}\label{dataset-and-classifiers}

Experiments use the ten-class FastAI Imagenette-v2 validation split,
containing 3,925 images. The four evaluated classifiers are ResNet-50,
ConvNeXt-Tiny, ViT-Base/16, and Swin-Tiny. All models remain frozen
during evaluation; the experimental variation is applied to their input
images.

Each regime evaluates every combination of six operator orders, five
severity levels, four models, and 3,925 images:

\[
6\times5\times4\times3{,}925=471{,}000
\]

compound evaluation records per regime. The clean baseline contributes
an additional 15,700 model-image records per regime. The resulting
accounting is 486,700 records per regime and 973,400 records across the
two regimes.

\subsection{3.2 Operator definitions and sequence
labels}\label{operator-definitions-and-sequence-labels}

The letters N, B, and J denote Gaussian noise, Gaussian blur, and JPEG
compression. Codes list operators in the order in which they are
applied, from left to right.

\begin{center}
\footnotesize
\begin{tabularx}{\columnwidth}{@{}lXl@{}}
\toprule
Code & Processing sequence & Final operator \\
\midrule
NBJ & Noise $\rightarrow$ Blur $\rightarrow$ JPEG & JPEG \\
NJB & Noise $\rightarrow$ JPEG $\rightarrow$ Blur & Blur \\
BNJ & Blur $\rightarrow$ Noise $\rightarrow$ JPEG & JPEG \\
BJN & Blur $\rightarrow$ JPEG $\rightarrow$ Noise & Noise \\
JNB & JPEG $\rightarrow$ Noise $\rightarrow$ Blur & Blur \\
JBN & JPEG $\rightarrow$ Blur $\rightarrow$ Noise & Noise \\
\bottomrule
\end{tabularx}
\end{center}

The nominal regime uses common severity parameters across orders and
deterministically seed-matched Gaussian-noise realizations. Image
identity, label, and model weights remain constant across paired order
comparisons.

\subsection{3.3 PSNR/SSIM-targeted parameter
regime}\label{psnrssim-targeted-parameter-regime}

The second regime selects sequence-specific parameter combinations using
a 50-image calibration subset, with five target bins defined by peak
signal-to-noise ratio (PSNR) and structural similarity (SSIM). We refer
to it as the \textbf{PSNR/SSIM-targeted regime}. These targets guided
parameter selection but were not achieved within the prespecified
tolerances, as discussed in Section 7.

\subsection{3.4 Outcomes and paired transition
measures}\label{outcomes-and-paired-transition-measures}

Top-1 accuracy is the primary outcome. PSNR and SSIM summarize low-level
similarity between each transformed image and its clean source. For each
model and sequence, accuracy is averaged across the five severity
settings and the 3,925 validation images. For each order, the
across-model mean is the arithmetic mean of the four model-specific
accuracies. Its six-order range is the maximum minus minimum of those
six means: 18.43 percentage points under nominal settings and 6.79 under
the PSNR/SSIM-targeted regime. These quantities are not the arithmetic
means of the four architecture-specific ranges (19.88 and 8.07 points,
respectively). Similarly, the 63.18\% contraction of the across-model
range differs from the mean of the four architecture-specific
contractions (58.08\%). Reported ranges are calculated from unrounded
accuracy values.

For two sequences A and B, let \(\hat{Y}_A\) and \(\hat{Y}_B\) be the
predicted classes and let \(C_A=\mathbb{1}[\hat{Y}_A=Y]\),
\(C_B=\mathbb{1}[\hat{Y}_B=Y]\) indicate correctness relative to the
true label \(Y\). Let \(L\) denote correct-to-incorrect transitions,
\(G\) incorrect-to-correct transitions, and \(M\) examples that are
incorrect under both sequences but receive different incorrect labels.
Then

\[
\begin{aligned}
\text{Correctness-state churn} &= P(C_A\ne C_B)=L+G,\\
\text{Predicted-label churn} &= P(\hat{Y}_A\ne\hat{Y}_B)=L+G+M.
\end{aligned}
\]

The five mutually exclusive states are: both correct; A correct/B
incorrect; A incorrect/B correct; both incorrect with the same predicted
class; and both incorrect with different predicted classes. For each
regime and model, the analysis compares all 15 unordered pairs of the
six processing orders. Each order-pair comparison aligns the same
validation image and severity index, giving 3,925 × 5 = 19,625 paired
image-severity observations per model and order pair. The regime-level
percentages pool these paired observations across the four models and
all 15 order pairs. The pooled unit is therefore an
image-severity-model-order-pair observation, not a unique image. All
transition measures are descriptive summaries of paired predictions.

\section{4. Results}\label{results}

\subsection{4.1 Order-specific
accuracy}\label{order-specific-accuracy}

The nominal results show substantial sensitivity to processing order
across all four architectures (Table 1). The model-specific six-order
ranges extend from 17.13 to 21.55 percentage points. ResNet-50 has the
largest range (21.55 pp), while Swin-Tiny has the smallest (17.13 pp).

\begin{table*}[!t]
\centering
\caption{Top-1 accuracy by processing order under nominal parameters (\%).}
\label{tab:nominal-accuracy}
\small
\setlength{\tabcolsep}{7pt}
\begin{tabular}{@{}lrrrrrrr@{}}
\toprule
Model & NBJ & NJB & BNJ & BJN & JNB & JBN & Range (pp) \\
\midrule
ResNet-50 & 53.44 & 56.88 & 37.30 & 35.71 & 57.27 & 35.79 & 21.55 \\
ConvNeXt & 63.98 & 66.44 & 48.63 & 45.99 & 65.44 & 46.33 & 20.45 \\
ViT-B/16 & 58.39 & 60.14 & 39.77 & 41.64 & 59.19 & 42.23 & 20.37 \\
Swin-T & 59.81 & 61.78 & 45.81 & 50.24 & 62.94 & 50.63 & 17.13 \\
\textbf{Mean} & \textbf{58.90} & \textbf{61.31} & \textbf{42.88} & \textbf{43.40} & \textbf{61.21} & \textbf{43.75} & \textbf{18.43} \\
\bottomrule
\end{tabular}
\vspace{3pt}
\begin{minipage}{\textwidth}
\footnotesize\emph{Note.} For each model, the range is the maximum minus minimum of its six order accuracies. In the final row, the six accuracies are four-model means and 18.43 is the range across those six means, not the average of the four model-specific ranges (19.88 pp). The ranges are calculated from unrounded values.
\end{minipage}
\end{table*}



\begin{figure*}[!t]
\centering
\includegraphics[width=0.88\textwidth]{RobustLens_Research_Figures/figure1_order_accuracy.png}
\caption{Order-specific accuracy for the nominal and PSNR/SSIM-targeted regimes.}
\label{fig:figure1-order-accuracy}
\end{figure*}

The final operator alone does not determine accuracy. NBJ and BNJ both
end in JPEG compression, yet their across-model mean accuracies are
58.90\% and 42.88\%, a difference of about 16 percentage points. The
operations applied before JPEG therefore matter for the image ultimately
presented to the classifier.

\subsection{4.2 Joint patterns in accuracy and image
fidelity}\label{joint-patterns-in-accuracy-and-image-fidelity}

Across the nominal orders, accuracy broadly follows the pattern of the
PSNR and SSIM summaries (Figure 2). Blur-ending orders NJB and JNB have
PSNR around 24.5-24.6 dB and SSIM around 0.68-0.69, while noise-ending
orders BJN and JBN have PSNR around 18.3 dB and SSIM around 0.27. There
is a useful exception: BNJ and BJN have similar mean accuracy (42.88\%
and 43.40\%) but different fidelity values (21.13 versus about 18.3 dB
PSNR; 0.386 versus about 0.27 SSIM). This is why accuracy and
image-fidelity metrics are best reported together. Because the
calibration tolerances were not met, these patterns are descriptive and
do not isolate a causal effect of fidelity.

\begin{figure*}[!t]
\centering
\includegraphics[width=0.68\textwidth]{RobustLens_Research_Figures/figure2_accuracy_fidelity.png}
\caption{Accuracy, PSNR, and SSIM across nominal corruption orders.}
\label{fig:figure2-accuracy-fidelity}
\end{figure*}

The nominal setting describes the end-to-end result of a specified
processing chain. The PSNR/SSIM-targeted regime asks how the results
change when sequence-specific parameters are selected toward shared
image-quality targets.

\subsection{4.3 Accuracy-range contraction under the
PSNR/SSIM-targeted
regime}\label{accuracy-range-contraction-under-the-psnrssim-targeted-regime}

Under the PSNR/SSIM-targeted regime, the range across the six
across-model mean accuracies decreases from 18.43 to 6.79 percentage
points, a 63.18\% descriptive contraction (Table 3). By architecture,
the largest contractions occur for ConvNeXt-Tiny (77.15\%) and ResNet-50
(72.62\%); ViT-B/16 contracts by 50.75\%, and Swin-Tiny by 31.80\%.

\begin{table*}[!t]
\centering
\caption{Top-1 accuracy by order under PSNR/SSIM-targeted parameters (\%).}
\label{tab:target-accuracy}
\small
\setlength{\tabcolsep}{7pt}
\begin{tabular}{@{}lrrrrrrr@{}}
\toprule
Model & NBJ & NJB & BNJ & BJN & JNB & JBN & Range (pp) \\
\midrule
ResNet-50 & 50.15 & 46.76 & 45.18 & 50.63 & 47.05 & 51.08 & 5.90 \\
ConvNeXt & 60.65 & 56.46 & 55.97 & 59.22 & 56.47 & 59.86 & 4.67 \\
ViT-B/16 & 56.54 & 53.39 & 46.51 & 54.72 & 53.01 & 55.68 & 10.03 \\
Swin-T & 57.61 & 50.44 & 53.93 & 61.23 & 52.31 & 62.12 & 11.68 \\
\textbf{Mean} & \textbf{56.23} & \textbf{51.76} & \textbf{50.40} & \textbf{56.45} & \textbf{52.21} & \textbf{57.18} & \textbf{6.79} \\
\bottomrule
\end{tabular}
\end{table*}

\begin{table*}[!t]
\centering
\caption{Change in six-order accuracy ranges.}
\label{tab:range-change}
\footnotesize
\begin{tabularx}{\textwidth}{@{}>{\raggedright\arraybackslash}X>{\centering\arraybackslash}p{1.35in}>{\centering\arraybackslash}p{1.55in}>{\centering\arraybackslash}p{1.45in}@{}}
\toprule
Architecture / summary & Nominal range (pp) & PSNR/SSIM-targeted range (pp) & Descriptive contraction (\%) \\
\midrule
ResNet-50 & 21.55 & 5.90 & 72.62 \\
ConvNeXt & 20.45 & 4.67 & 77.15 \\
ViT-B/16 & 20.37 & 10.03 & 50.75 \\
Swin-T & 17.13 & 11.68 & 31.80 \\
\textbf{Across-model range} & \textbf{18.43} & \textbf{6.79} & \textbf{63.18} \\
\bottomrule
\end{tabularx}
\vspace{3pt}
\begin{minipage}{\textwidth}
\footnotesize\emph{Note.} The final row reports the range of the six across-model mean accuracies in each regime; it is not the mean of the four architecture-specific ranges (19.88 pp nominal; 8.07 pp PSNR/SSIM-targeted). The 63.18\% contraction of this aggregate range differs from the arithmetic mean of the four architecture-specific contractions (58.08\%).
\end{minipage}
\end{table*}



\begin{figure*}[!t]
\centering
\includegraphics[width=0.80\textwidth]{RobustLens_Research_Figures/figure3_range_contraction.png}
\caption{Observed six-order accuracy ranges in the nominal and PSNR/SSIM-targeted regimes.}
\label{fig:figure3-range-contraction}
\end{figure*}

The results show why benchmark reports should state both the
parameter-selection protocol and the measured image-fidelity profile.
The contraction describes the difference between these two tested
regimes; it is not an estimate under exact equal-fidelity conditions.

\subsection{4.4 Paired predictions expose changes beyond
correctness}\label{paired-predictions-expose-changes-beyond-correctness}

Aggregate accuracy does not reveal whether two sequences produce the
same predictions on the same images. The five-state paired analysis
separates changes in correctness from changes in the identity of
incorrect predictions. Under nominal parameters, correctness-state churn
is 16.93\%, while predicted-label churn is 40.78\%. Under the
PSNR/SSIM-targeted regime, the corresponding values are 12.34\% and
34.74\% (Table 4).

\begin{table*}[!t]
\centering
\caption{Paired-prediction transition summaries (\%).}
\label{tab:prediction-churn}
\small
\begin{tabularx}{\textwidth}{@{}>{\raggedright\arraybackslash}Xrr@{}}
\toprule
Measure & Nominal parameters & PSNR/SSIM-targeted regime \\
\midrule
Correctness-state churn, $L+G$ & 16.93 & 12.34 \\
Error migration, $M$ & 23.85 & 22.40 \\
\textbf{Predicted-label churn, $L+G+M$} & \textbf{40.78} & \textbf{34.74} \\
\bottomrule
\end{tabularx}
\end{table*}

\begin{figure*}[!t]
\centering
\includegraphics[width=0.90\textwidth]{RobustLens_Research_Figures/figure4_prediction_churn.png}
\caption{Predicted-label churn decomposed into correctness-state churn and migration between incorrect labels.}
\label{fig:figure4-prediction-churn}
\end{figure*}

The component \(M\) counts cases that remain incorrect but receive a
different incorrect label under the second sequence. It accounts for
23.85\% of paired observations under nominal settings and 22.40\% under
the PSNR/SSIM-targeted regime. These transitions are invisible to
correctness-state churn, even though the model output has changed.

The class-level results show the same pattern for ``church'': nominal
accuracy is 18.35\%, correctness-state churn is 11.04\%, error migration
is 39.41\%, and predicted-label churn is 50.45\%. This class-specific
result is descriptive and is reported alongside the aggregate analysis.

\section{5. Discussion}\label{discussion}

Processing order materially changes classification accuracy in this
experiment. With fixed images and model weights, the six-order accuracy
ranges are 17.13-21.55 percentage points across architectures; even
orders with the same final operator can differ by about 16 points.
Reporting the full chain therefore preserves information that a label
based only on the last corruption would miss.

The second result is that the range of across-model mean accuracy is
63.18\% smaller under the PSNR/SSIM-targeted regime. The
architecture-specific contractions vary: ConvNeXt-Tiny and ResNet-50
contract most, whereas ViT-B/16 and Swin-Tiny contract less. The
calibration tolerances were not met, so this comparison describes the
observed response to parameter selection rather than isolating a pure
order effect at equal fidelity.

The paired analysis shows a separate change that aggregate accuracy
cannot reveal. Under nominal settings, predicted-label churn is 40.78\%,
compared with correctness-state churn of 16.93\%; 23.85\% of paired
observations remain incorrect while switching to a different incorrect
class. The same pattern remains under the PSNR/SSIM-targeted regime.
Together, accuracy, measured fidelity, and paired class transitions give
a fuller description of the evaluated pipelines than any one summary
alone.

\section{6. Conclusion}\label{conclusion}

This experiment compares all six orders of Gaussian noise, Gaussian
blur, and JPEG compression across four frozen vision classifiers. The
model-specific nominal six-order ranges span 17.13 to 21.55 percentage
points, and the range of across-model mean accuracy contracts by 63.18\%
under the PSNR/SSIM-targeted regime. Paired predictions also show that
nominal predicted-label churn (40.78\%) substantially exceeds
correctness-state churn (16.93\%), because some images switch between
different incorrect classes. These results make the case for reporting
order-specific accuracy and fidelity alongside paired prediction
transitions.

\section{7. Limitations and future work}\label{limitations-and-future-work}

The next step is to repeat the experiment on a larger label space,
additional datasets, and more model families, including external
validation on ImageNet-V2, and to report uncertainty intervals for the
range and transition summaries. The present results are limited to the
ten-class Imagenette-v2 validation split and four frozen architectures.
The PSNR/SSIM-targeted regime did not meet the prespecified calibration
tolerances (PSNR RMSE 0.943 dB versus 0.75 dB; SSIM RMSE 0.138 versus
0.035). Accordingly, the 63.18\% contraction is a descriptive difference
between the two tested parameter regimes, not an estimate under exact
equal-fidelity conditions.

\section*{References}\label{references}

{[}1{]} D. Hendrycks and T. Dietterich. ``Benchmarking Neural Network
Robustness to Common Corruptions and Perturbations.''
\emph{International Conference on Learning Representations (ICLR)},
2019. \href{https://arxiv.org/abs/1903.12261}{arXiv:1903.12261}

{[}2{]} B. Goyal, J.-F. Lalonde, Y. Li, and M. Gupta. ``Robust Scene
Inference under Noise-Blur Dual Corruptions.'' \emph{IEEE International
Conference on Computational Photography (ICCP)}, 2022.
\href{https://arxiv.org/abs/2207.11643}{arXiv:2207.11643}

{[}3{]} X. Liao, K. Li, X. Zhu, and K. J. R. Liu. ``Robust Detection of
Image Operator Chain With Two-Stream Convolutional Neural Network.''
\emph{IEEE Journal of Selected Topics in Signal Processing},
14(5):955-968, 2020.
\href{https://doi.org/10.1109/JSTSP.2020.3002391}{doi:10.1109/JSTSP.2020.3002391}

{[}4{]} M. Milani Fard, Q. Cormier, K. Canini, and M. R. Gupta. ``Launch
and Iterate: Reducing Prediction Churn.'' \emph{Advances in Neural
Information Processing Systems}, 29:3179-3187, 2016.
\href{https://proceedings.neurips.cc/paper_files/paper/2016/hash/dc5c768b5dc76a084531934b34601977-Abstract.html}{NeurIPS
proceedings}

{[}5{]} C. Li et al.~``R-Bench: Are Your Large Multimodal Model Robust
to Real-World Corruptions?'' \emph{IEEE Journal of Selected Topics in
Signal Processing}, 19(7):1349-1361, 2025.
\href{https://doi.org/10.1109/JSTSP.2025.3558652}{doi:10.1109/JSTSP.2025.3558652}

{[}6{]} O. Dünkel, A. Jesslen, J. Xie, C. Theobalt, C. Rupprecht, and A.
Kortylewski. ``CNS-Bench: Benchmarking Image Classifier Robustness Under
Continuous Nuisance Shifts.'' \emph{Proceedings of the IEEE/CVF
International Conference on Computer Vision (ICCV)}, 19978-19988, 2025.
\href{https://openaccess.thecvf.com/content/ICCV2025/html/Dunkel_CNS-Bench_Benchmarking_Image_Classifier_Robustness_Under_Continuous_Nuisance_Shifts_ICCV_2025_paper.html}{CVF
Open Access}

{[}7{]} F. Li, T. Klein, W. Brendel, R. Geirhos, and R. S. Zimmermann.
``LAION-C: An Out-of-Distribution Benchmark for Web-Scale Vision
Models.'' \emph{Proceedings of the 42nd International Conference on
Machine Learning}, PMLR 267:35136-35167, 2025.
\href{https://proceedings.mlr.press/v267/li25aw.html}{PMLR paper page}

{[}8{]} Y. Wei, C. Zhang, Y. Wang, X. Wang, B. Shao, H. Li, and L. Di.
``What Paired Evaluations Reveal under Visual Perturbations.''
arXiv:2609.35583, 2026.
\href{https://arxiv.org/abs/2609.35583}{arXiv:2609.35583}

{[}9{]} I. Mason, A. Sarkar, T. Sasaki, and X. Boix. ``Modularity Trumps
Invariance for Compositional Robustness.'' arXiv:2306.09005, 2023.
\href{https://arxiv.org/abs/2306.09005}{arXiv:2306.09005}

\end{document}
